% PSEUDOSLICE1
\documentclass[11pt]{article}
\usepackage[margin=.88in]{geometry}
\usepackage{amsmath,amssymb,amsthm}
\newtheorem{theorem}{Theorem}
\newtheorem{proposition}{Proposition}
\begin{document}

\begin{center}
{\Large PSEUDOSLICE1: Exact-slice pseudo-independence under the 2026 method-of-moments definition}\\[4pt]
\textbf{Date: 2026-10-04}\\
\textbf{Status: PROVED/EXACT reductions and lower bound + COMPUTATIONAL LP}\\
\textbf{P versus NP status: OPEN}
\end{center}

\section{Question}

Cavalar--Fabris--Mukhopadhyay--Srinivasan--Yehudayoff (ECCC TR26-220, 2026)
prove a robust sunflower theorem for distributions that are locally
pseudo-independent. Their Definition 2.2 asks that every marginal on at most
\(k\) coordinates be a convex combination of product distributions whose
coordinate biases all lie in
\[
[p(1-\varepsilon),p(1+\varepsilon)].
\]

For our negative distribution
\[
\mu_- = U_{M,t},
\]
the uniform exact-\(t\) slice on \(M\) clause coordinates, the question is:
how small can \(\varepsilon\) be for local order \(k\)?

\section{Exact LP reduction}

Let
\[
p=\frac tM.
\]
For a fixed set of \(k\) coordinates, the marginal of \(U_{M,t}\) is
exchangeable.

\begin{theorem}[SLICE-PSEUDO-LP]
For fixed \(M,t,k,\varepsilon\), membership of the \(k\)-coordinate marginal
of \(U_{M,t}\) in the pseudo-independent class of Definition 2.2 is exactly
equivalent to feasibility of a linear program with only \(k+1\) nonnegative
mixture weights.
\end{theorem}

\begin{proof}
Any product distribution with coordinate probabilities
\(p_i\in[p(1-\varepsilon),p(1+\varepsilon)]\) is multi-affine in the
\(p_i\)'s. Hence it lies in the convex hull of the \(2^k\) corner product
distributions with
\[
p_i\in\{p(1-\varepsilon),p(1+\varepsilon)\}.
\]
Since the target marginal is invariant under all coordinate permutations, any
feasible mixture may be symmetrized. The \(2^k\) corners therefore collapse
to \(k+1\) types, indexed by the number \(j\) of high-bias coordinates.

Exchangeable laws on \(\{0,1\}^k\) are determined by their factorial moments.
For \(r=0,\ldots,k\), the target normalized moment is
\[
m_r
=
\frac{(t)_r/(M)_r}{p^r}.
\]
For the symmetrized corner type with \(j\) high coordinates, the normalized
\(r\)-moment equals
\[
A_{rj}(\varepsilon)
=
\frac1{\binom kr}
\sum_\ell
\binom j\ell
\binom{k-j}{r-\ell}
(1+\varepsilon)^\ell
(1-\varepsilon)^{r-\ell}.
\]
Thus feasibility is exactly
\[
Aw=m,
\qquad
w_j\ge0,
\qquad
\sum_j w_j=1.
\]
Projection of a feasible \(k\)-coordinate representation gives all smaller
marginals, so checking order \(k\) suffices.
\end{proof}

\section{Necessary epsilon from the pair deficit}

Define
\[
d
=
1-
\frac{(t)_2/(M)_2}{p^2}
=
\frac{M-t}{t(M-1)}.
\]

\begin{theorem}[SLICE-PSEUDO-LOWER]
If \(U_{M,t}\) is \((k,\varepsilon)\)-pseudo-independent in the sense of
TR26-220, then
\[
\boxed{
\varepsilon^2
\ge
(k-1)\frac{M-t}{t(M-1)}
}.
\]
For \(k=2\), equality is achievable.
\end{theorem}

\begin{proof}
Symmetrize a feasible mixture and write each component bias as
\[
p_i=p(1+u_i),
\qquad |u_i|\le\varepsilon.
\]
The exact one-coordinate marginal gives
\[
\mathbb E[u_1]=0.
\]
The exact two-coordinate inclusion moment gives
\[
\mathbb E[u_1u_2]=-d.
\]
Therefore
\[
0
\le
\mathbb E\left(\sum_{i=1}^k u_i\right)^2
=
k\mathbb E[u_1^2]
+
k(k-1)\mathbb E[u_1u_2].
\]
Hence
\[
\mathbb E[u_1^2]\ge(k-1)d.
\]
But \(u_1^2\le\varepsilon^2\), proving the bound.

For \(k=2\), mix equally the product distributions with normalized bias
vectors \((+\varepsilon,-\varepsilon)\) and
\((-\varepsilon,+\varepsilon)\), with \(\varepsilon^2=d\).
This matches the exact one- and two-coordinate moments.
\end{proof}

\section{Consequence for Theorem 2.3 parameters}

TR26-220 requires
\[
\varepsilon\le\frac\gamma w
\]
for a sufficiently small absolute constant \(\gamma\), where \(w\) is the
DNF width. Combining this with SLICE-PSEUDO-LOWER yields the necessary
condition
\[
k-1
\le
\frac{\gamma^2}{w^2}
\frac{t(M-1)}{M-t}.
\tag{1}
\]

Their sunflower corollary uses local order
\[
k=\Theta\left(w\log\frac1\delta\right).
\]
Therefore a direct application to an exact sparse slice can only possibly
operate in the asymptotic window
\[
\boxed{
w^3\log(1/\delta)=O(t)
}
\tag{2}
\]
up to the theorem constants.

Equation (2) is a necessary condition, not a sufficient one.

For our SAT slice \(t=\Theta(q)\), this places the potential width scale at
no larger than roughly \(q^{1/3}\) up to logarithms.

\section{Numerical exact-LP experiment}

The exact \((k+1)\)-variable LP above was solved using normalized factorial
moments. Each entire \(q\) run took under one second.

For stability we report \(k\le6\).

\[
\begin{array}{c|ccccc}
q & \varepsilon_2 & \varepsilon_3 & \varepsilon_4 &
\varepsilon_5 & \varepsilon_6 \\\
\hline
16  & .1328 & .2356 & .3238 & .4028 & .4754 \\
20  & .1191 & .2108 & .2891 & .3591 & .4233 \\
24  & .1088 & .1923 & .2634 & .3268 & .3849 \\
32  & .0944 & .1663 & .2273 & .2816 & .3311 \\
64  & .0668 & .1171 & .1595 & .1969 & .2309 \\
128 & .0472 & .0826 & .1121 & .1380 & .1615 \\
256 & .0334 & .0582 & .0789 & .0970 & .1133
\end{array}
\]

For \(k=2\), the LP agrees with the exact formula
\[
\varepsilon_2
=
\sqrt{
\frac{M-t}{t(M-1)}
}
\]
to numerical precision.

For fixed small \(k\), the data are consistent with
\[
\varepsilon_k=\Theta_k(t^{-1/2}),
\]
but no uniform-in-\(k\) asymptotic upper bound is claimed.

\section{Interpretation}

This changes the status of the 2026 pseudo-independence theorem.

It is \textbf{not} automatically incompatible with an exact Hamming slice:
the slice marginal can indeed be represented locally as a convex mixture of
nearby product distributions.

At the same time, the fixed-weight negative dependence imposes a quantitative
cost, captured by SLICE-PSEUDO-LOWER. The theorem can only help in a bounded
width/local-order regime satisfying (1).

The remaining larger problem is the positive distribution. Pair balancing
gives exact degree-two orbit moments, but TR26-220 needs pseudo-independence
through local order \(k=\Theta(w\log(1/\delta))\). Establishing such a
property for an UNSAT-supported distribution is now a concrete target.

\section{Next target}

\textbf{POS-PSEUDO3.}
For the pair-balanced positive UNSAT distribution, estimate every
three-coordinate local marginal from exact canonical triple-orbit counts and
compute the minimum product-box \(\varepsilon\) needed for each orbit.
This is an exact 8-corner convex-hull test once the triple inclusion
probability is specified.

If the worst canonical triple requires very large \(\varepsilon\), the
pseudo-independence bridge fails already at order 3. If all triple orbits are
close to the exact-slice value, proceed to order 4 in bounded chunks.

Mandatory guards remain large-core structure and XOR-SLICE-EASY.

\section{Status}

SLICE-PSEUDO-LP: \textbf{PROVED/EXACT}.

SLICE-PSEUDO-LOWER: \textbf{PROVED/EXACT}.

Finite LP values: \textbf{COMPUTATIONAL}.

Useful higher-order pseudo-independence of the positive UNSAT distribution:
\textbf{OPEN}.

TWO-DIST-MOMENT-APPROX and the monotone slice lower bound:
\textbf{OPEN}.

Therefore \(P\ne NP\) remains OPEN.

\begin{thebibliography}{3}
\bibitem{CFMSY2026}
B. P. Cavalar, T. Fabris, P. Mukhopadhyay, S. Srinivasan, and A. Yehudayoff,
Derandomized Sunflowers and Radical Lower Bounds,
ECCC TR26-220, 2026.

\bibitem{CKR2022}
B. P. Cavalar, M. Kumar, and B. Rossman,
Monotone Circuit Lower Bounds from Robust Sunflowers,
Algorithmica 84:3655--3685, 2022.

\bibitem{Wegener1985}
I. Wegener,
On the Complexity of Slice Functions,
Theoretical Computer Science 38:55--68, 1985.
\end{thebibliography}

\end{document}
